\documentclass[12pt]{article}

% Page setup
\usepackage[margin=1in]{geometry}
\usepackage{setspace}
\usepackage{enumitem}
\usepackage{amsmath, amssymb}

% No paragraph indent
\setlength{\parindent}{0pt}

% Line spacing
\renewcommand{\baselinestretch}{1.2}

\begin{document}

\begin{center}
\textbf{\Large MATH 2790 Fall 2025}\\[1em]
\textbf{\Large Homework 7}\\[0.5em]
Due date: Wednesday, Oct 15
\end{center}

\bigskip

\textbf{Instructions:}

\begin{enumerate}[label=(\alph*)]
\item It is expected to show detailed work for every problem. You are allowed to have a
discussion with peers on some challenging questions. But then you should work on
your own to finish the assignment.

\item All questions in this assignment are from sections 3.4--3.6.
\end{enumerate}

\bigskip

\textbf{Problems:}

\begin{enumerate}
\item Question (1) on page 183 of textbook under section 3.4.

\item Question (2) on page 184 of textbook under section 3.4.

\item Question (2) on page 185 of textbook under section 3.5.

\item Question (3) on page 185 of textbook under section 3.5.

\item Question (1) on page 186 of textbook under section 3.6.

\item Question (2) on page 186 of textbook under section 3.6.

\item Question (5) on page 186 of textbook under section 3.6.

\item A loan is to be repaid at a nominal interest rate of 6\% payable monthly for 60
months. The monthly payments of 1000 are paid at the end of each month. Use two
methods to calculate the outstanding loan balance right before the 40th payment.
\end{enumerate}

\end{document}
